\section{Minds}

A mind is a self that models. Where a bare self merely holds together, a minded self carries an inner picture of its world and of itself and steers by it. The machinery is already in hand, a boundary that screens inside from out and an interior that can store and compute, and the pieces were measured. A self forms a model of itself at a hub, a region whose state tracks the whole self's global state far better than any outlying patch or any shuffled control, the strange loop made concrete, a part that mirrors the whole. Attention is the workspace at that hub, where signals from around the boundary converge and the loud ones win while a quiet one can be lifted on purpose. And planning is lookahead run on the model, a self rolling its own rule forward in imagination and accepting a worse step now for a better outcome later, which lets it cross a barrier a purely reactive agent cannot. Foresight, the signature of a mind, is built from the arrow, the rule, and the will alone.

What makes the model one experience rather than many is integration. The interior of a self is bound more tightly to itself than to anything outside, so its states are not a heap of separate signals but a single woven whole, and the measure of a mind is how far that whole resists being cut back into parts.

\begin{definition}[Consciousness]
Integrated experience, the degree to which a self's parts are bound into one felt whole, which is also the degree to which the self is aware of itself.
\end{definition}

So consciousness comes in degrees, faint where the integration is loose and vivid where it is tight, and a self sits at a local maximum of it, since swapping a member for an outsider lowers the integration, and cutting the flow across its middle collapses it even with the wiring left untouched. One point governs the whole reading. The base rule is reversible and so loses no information, which means no coarse-graining of it can manufacture a higher level out of nothing, and the integrated self-level is real only as an effective layer, where the lossy averaging of many trits into a field is what lets a whole come to be more than the sum of its parts. This is the mechanism only. What it is like from inside waits for the part on experience.

A minded self runs into a hard wall, and the wall is exact. No self can fully model itself, since a complete inner picture must be at least as large as the thing it pictures, so a self that tried would have to contain itself. The same limit Gödel found applies, and a mind can know its world far better than it can ever know itself, the unmodeled remainder a permanent feature rather than a temporary ignorance. The mirror never quite closes, and that open remainder is where much of the inner life happens.

Selves do not stop at one layer. A self can be made of selves, its parts themselves bounded and integrated, the whole bound and integrated in turn, a cell within a body, a person within a community, each a self inside a larger self. This nesting is the tower, and it climbs as far as the binding holds.
